\subsection{Solve Any Equation} 
In working with radical equations, we have two new issues in solving equations. \newstuff{New moves in red, below.}
\begin{fact}[Solve Any Equation]
\textbf{Basic moves}
\begin{itemize}
\item Simplify the expressions on either side.
\item Add or subtract the same expression on both sides of the equation.
\item Multiply or divide by the same expression on both sides of the equation.
\begin{itemize}
\item We must be careful not to multiply or divide by zero.
\item If we accidentally multiply by zero, we introduce extraneous solutions.
\item If we multiplied by a variable expression, we must check.
\end{itemize}
\end{itemize}
\textbf{How does the variable appear?}
\begin{itemize}
\item Once, inside a function {\textrightarrow} Isolate, then eliminate the function
\begin{itemize}
\item For a one-to-one function $f$ and $c$ in the range of $f$, $f(X)=c$ and $X=f\inv(c)$ are equivalent.
\item If the function is not one-to-one, there may be many solutions: we cut up the domain of $f$ to have a partial inverse.
\item If the function's range is not all real numbers, there may be no solutions solutions.
\item \newstuff{If the \emph{inverse} is not one-to-one, we may introduce \term{extraneous solutions}, and must check.}
\end{itemize}
\item \newstuff{Many times, inside a function {\textrightarrow} Isolate, then eliminate one instance of the function.}
\item Many times, as a polynomial {\textrightarrow} Write in standard form, factor, use the zero property.
\end{itemize}
\end{fact}